\documentclass[conference]{IEEEtran}
\IEEEoverridecommandlockouts
\usepackage[utf8]{inputenc}
\usepackage[T1]{fontenc}
\usepackage[spanish,es-noshorthands,es-tabla,es-nodecimaldot]{babel}
\usepackage{amsmath,amssymb,amsfonts}
\usepackage{graphicx}
\usepackage{booktabs}
\usepackage{xcolor}
\usepackage[hyphens]{url}
\usepackage{cite}
\usepackage{tikz}
\usetikzlibrary{positioning,arrows.meta,fit}
\usepackage[hidelinks]{hyperref}

% ---------------------------------------------------------------------------
% Marcador de texto pendiente: se completa con los números de la corrida `full`.
% Antes de entregar, buscar "\pend" y reemplazar cada uno (o dejar \pend vacío).
\newcommand{\pend}[1]{\textcolor{red}{[#1]}}
% Números generados automáticamente por paper/make_tables.py (si existen)
\providecommand{\HW}{\pend{GPU}}\providecommand{\GPUH}{\pend{h}}
\providecommand{\NTRAIN}{\pend{n}}\providecommand{\NSYN}{\pend{n}}
\providecommand{\NVAL}{\pend{n}}\providecommand{\NTEST}{\pend{n}}
\providecommand{\VTORCH}{\pend{v}}\providecommand{\VTRANSF}{\pend{v}}
\InputIfFileExists{generated/numeros.tex}{}{}
% Tabla o figura generada; si aún no existe, muestra un recuadro pendiente.
\newcommand{\tablagen}[1]{\InputIfFileExists{generated/#1.tex}{}{\fbox{\pend{\detokenize{#1}: ejecutar paper/make\_tables.py}}}}
\newcommand{\figgen}[2]{\IfFileExists{figs/#1.png}{\includegraphics[width=#2]{figs/#1.png}}{\fbox{\parbox{0.9#2}{\centering\pend{\detokenize{#1}.png pendiente}}}}}
% ---------------------------------------------------------------------------

\begin{document}

\title{Diagnóstico de fallas en rodamientos con pocos datos de falla y cambio de carga: CNN, Audio Spectrogram Transformer y aumento por difusión condicional}

\author{
\IEEEauthorblockN{Hipólito Coaguila Berrios, Lenin Dixon Herrera Padilla, Jhafet Cánepa Maceda}
\IEEEauthorblockA{Maestría en Inteligencia Artificial, Universidad Nacional de Ingeniería (UNI), Lima, Perú\\
MIA-203 Redes Neuronales y Aprendizaje Profundo --- Grupo 6\\
Código: \url{https://github.com/jafcn09/examenfinal_redes_uniingenieria}}
}

\maketitle

\begin{abstract}
Los rodamientos son la principal causa de falla de los motores eléctricos y su diagnóstico automático a partir de la vibración enfrenta dos limitaciones reales: hay muy pocos registros de fallas, sobre todo de las severas, y la máquina opera en condiciones distintas de las vistas al entrenar. Estudiamos la clasificación de 10 condiciones de rodamiento (normal y tres tipos de falla con tres severidades) del conjunto CWRU bajo un escenario que reproduce ambas limitaciones: entrenamiento con cargas de 0 a 2~HP con desbalance inducido (30\,\%, 10\,\% y 3\,\% de las ventanas según severidad) y prueba en la carga de 3~HP nunca vista. Comparamos tres arquitecturas ---una CNN 1D sobre la señal cruda (línea base), una CNN 2D sobre el espectrograma y un Audio Spectrogram Transformer (AST) preentrenado en AudioSet--- bajo cuatro estrategias de aumento: ninguna, sobremuestreo de las ventanas reales (control), clásica y sintética, generada por un único modelo de difusión condicional 1D (DDPM con guía sin clasificador). Con 3 semillas, el mejor sistema alcanza un F1-macro de \pend{X.XXX $\pm$ X.XXX} en 3~HP (\pend{modelo + aumento}), frente a \pend{X.XXX} de la línea base; con ruido de 0~dB, \pend{resultado}. El análisis Grad-CAM muestra que \pend{\% de las ventanas} concentran la activación en la periodicidad de las frecuencias teóricas de falla (BPFO/BPFI). El costo total fue de \GPUH~horas de GPU. Concluimos que \pend{conclusión principal en una oración}.
\end{abstract}

\begin{IEEEkeywords}
diagnóstico de fallas, rodamientos, modelos de difusión, Audio Spectrogram Transformer, aumento de datos, desbalance de clases, Grad-CAM.
\end{IEEEkeywords}

% ===========================================================================
\section{Introducción}
Los rodamientos concentran el 41--44\,\% de las fallas de los motores eléctricos industriales~\cite{ieee1985motor}. Una falla no detectada detiene la máquina y todo lo que depende de ella; en la industria peruana esto alcanza a molinos, fajas y bombas de la minería, a las estaciones de bombeo de agua y a las bombas de riego y cámaras de frío de la agroindustria. El \emph{mantenimiento predictivo} busca detectar el daño a tiempo a partir de un sensor de vibración y decidir cuándo intervenir.

El aprendizaje profundo aprende directamente de la señal las características que antes se diseñaban a mano (envolvente, curtosis, bandas espectrales)~\cite{randall2011tutorial,zhao2020benchmark}, pero en planta enfrenta dos obstáculos que el uso habitual del conjunto CWRU oculta~\cite{smith2015cwru,neupane2020review}: (i) los datos de falla son escasos ---casi todo el registro corresponde a la máquina sana y las fallas severas son las más raras--- y (ii) la condición de operación cambia (carga, velocidad, ruido). Este trabajo formula el diagnóstico en un escenario que reproduce ambas dificultades y evalúa si la generación de señales sintéticas con un modelo de difusión mitiga la escasez de fallas, y si arquitecturas modernas como el Transformer preentrenado aportan sobre una CNN.

\textbf{Contribuciones.}
\begin{enumerate}
\item Un protocolo de evaluación exigente sobre CWRU: partición por carga (0--2~HP para entrenar, 3~HP para probar), desbalance de cola larga por severidad, evaluación con ruido (10 y 0~dB) y 3 semillas, sin fuga de datos entre particiones.
\item Una comparación controlada $3\times4$ entre arquitecturas (CNN 1D, CNN 2D y AST) y estrategias de aumento (ninguna, sobremuestreo, clásica y difusión condicional con un único generador para las 10 clases); el sobremuestreo separa el efecto de equilibrar las clases del aporte de las muestras generadas.
\item Un análisis de interpretabilidad que contrasta Grad-CAM y la atención del AST con las frecuencias de falla que predice la cinemática del rodamiento (BPFO, BPFI), y una evaluación de la calidad de las señales sintéticas (t-SNE y TSTR).
\item La corrección documentada de dos problemas del conjunto CWRU (frecuencia de muestreo de la línea base normal y un archivo con una señal duplicada), y un repositorio reproducible.
\end{enumerate}

% ===========================================================================
\section{Trabajos relacionados}
\textbf{Diagnóstico de rodamientos y CWRU.} Randall y Antoni~\cite{randall2011tutorial} establecen el diagnóstico clásico por análisis de envolvente alrededor de las frecuencias características de falla. Smith y Randall~\cite{smith2015cwru} analizan el conjunto CWRU archivo por archivo y muestran que varias condiciones no son diagnosticables con claridad, por lo que exactitudes cercanas al 100\,\% deben leerse con cautela. Las revisiones de Neupane y Seok~\cite{neupane2020review} y de Zhao~\textit{et al.}~\cite{zhao2020benchmark} recopilan modelos de aprendizaje profundo sobre CWRU y advierten que muchas particiones aleatorias de ventanas solapadas sobreestiman el desempeño; nuestro protocolo separa por carga y en el tiempo para evitarlo. Paderborn~\cite{lessmeier2016kat} ofrece daños reales por desgaste y queda como validación futura.

\textbf{Redes convolucionales sobre la señal.} WDCNN~\cite{zhang2017wdcnn} propone una CNN 1D con un primer núcleo ancho que actúa como banco de filtros y reporta robustez al ruido y a cambios de carga; la usamos como línea base.

\textbf{Transformers para audio.} El Transformer~\cite{vaswani2017attention} y su versión para imágenes por parches, ViT~\cite{dosovitskiy2021vit}, dieron lugar al AST~\cite{gong2021ast}, que aplica atención sobre parches del espectrograma mel y se preentrena en AudioSet~\cite{gemmeke2017audioset}. A diferencia de los trabajos que entrenan Transformers desde cero sobre vibración, evaluamos si el preentrenamiento en audio general se transfiere a vibración industrial.

\textbf{Modelos generativos para el desbalance.} Los modelos de difusión (DDPM)~\cite{ho2020ddpm} generan muestras invirtiendo un proceso de ruido; DDIM~\cite{song2021ddim} acelera el muestreo y la guía sin clasificador~\cite{ho2022cfg} controla la fidelidad a la clase. Yang~\textit{et al.}~\cite{yang2024ddpmfault} y Wang~\textit{et al.}~\cite{wang2024ddimtransnet} aplican difusión al diagnóstico con datos desbalanceados, pero generan imágenes tiempo-frecuencia (CWT, GADF) y prueban en la misma condición de operación. Nosotros generamos directamente la señal 1D con un único modelo condicional, la evaluamos bajo cambio de carga y la comparamos contra aumento clásico~\cite{park2019specaugment}, que es el control necesario para atribuir la mejora al modelo generativo.

\textbf{Interpretabilidad.} Grad-CAM~\cite{selvaraju2017gradcam} y el \emph{rollout} de atención~\cite{abnar2020rollout} explican la decisión; la evaluación TSTR (entrenar con sintético, probar con real)~\cite{esteban2017rcgan} mide la utilidad de los datos generados.

% ===========================================================================
\section{Metodología}
\subsection{Datos}
Usamos el CWRU Bearing Data Center~\cite{cwru}: vibración del acelerómetro del lado de accionamiento (\emph{drive end}) de un motor de 2~HP con rodamiento SKF~6205, a 12~kHz, para cuatro cargas (0--3~HP, 1797--1730~rpm). Las 10 clases son: Normal y fallas en la bola (B), la pista interna (IR) y la pista externa en la posición de las 6 (OR), cada una con diámetros de 0,007, 0,014 y 0,021~in; en total, 40 archivos de unos 10~s. Se excluyen el daño de 0,028~in (no existe para OR), las posiciones OR @3 y @12 y los datos del lado del ventilador. La fuente es pública para uso académico; los archivos se descargan de un espejo con los \texttt{.mat} originales.

\textbf{Correcciones de datos.} Verificamos que la línea base Normal está registrada a 48~kHz y no a 12~kHz: sus líneas espectrales (159,5, 360 y 533,5~Hz) coinciden con las de los archivos de falla solo al leerla a 48~kHz, y su duración en muestras es cuatro veces mayor. Sin corrección, la clase Normal tiene otra escala temporal y espectral y la red puede reconocerla por ese artefacto; la remuestreamos a 12~kHz con un filtro polifásico. Además, el archivo \texttt{99\_Normal\_2.mat} contiene también la señal de \texttt{98} (carga 1); seleccionamos la variable que corresponde al número del archivo.

\textbf{Ventanas y partición.} Cada señal se divide en ventanas de $L=2048$ muestras (171~ms, unas 5 revoluciones) normalizadas por ventana ($z$-score). Entrenamiento: cargas 0--2~HP, el 80\,\% inicial de cada señal con paso 512; validación: el 20\,\% final con paso 2048 (sin solape con entrenamiento); prueba: la carga de 3~HP completa con paso 2048, nunca usada para elegir épocas ni hiperparámetros. El desbalance se induce solo en entrenamiento, conservando el 30\,\%, 10\,\% y 3\,\% de las ventanas de las severidades 0,007, 0,014 y 0,021~in. La Tabla~\ref{tab:particiones} resume las particiones (\NTRAIN{} ventanas de entrenamiento, \NVAL{} de validación y \NTEST{} de prueba).

\begin{table}[t]
\centering
\caption{Ventanas por clase y partición.}
\label{tab:particiones}
\footnotesize
\tablagen{tabla_particiones}
\end{table}

\textbf{Física de la falla.} Un defecto produce impactos periódicos a una frecuencia que depende de la geometría (número de bolas $N_b=9$, diámetro de bola $d$, diámetro primitivo $D$, ángulo de contacto $\phi$) y de la frecuencia de giro $f_r$~\cite{randall2011tutorial}:
\begin{align}
\mathrm{BPFO} &= \tfrac{N_b}{2} f_r \left(1-\tfrac{d}{D}\cos\phi\right) = 3{,}585\,f_r, \label{eq:bpfo}\\
\mathrm{BPFI} &= \tfrac{N_b}{2} f_r \left(1+\tfrac{d}{D}\cos\phi\right) = 5{,}415\,f_r, \label{eq:bpfi}
\end{align}
es decir, 103~Hz y 156~Hz a 1730~rpm. Estas frecuencias sirven para verificar qué aprendió la red (Sección~\ref{sec:interp}).

\subsection{Arquitecturas}
La Fig.~\ref{fig:pipeline} resume el flujo.

\begin{figure}[t]
\centering
\begin{tikzpicture}[node distance=2.5mm and 3mm, font=\scriptsize,
  box/.style={draw, rounded corners=1.5pt, align=center, minimum height=6mm, inner sep=2pt, fill=#1},
  box/.default=white, >={Stealth[length=1.6mm]}]
\node[box=gray!10] (sig) {Señal 12\,kHz\\ventana 2048};
\node[box=orange!15, below=of sig] (ddpm) {DDPM 1D condicional\\U-Net + FiLM, CFG};
\node[box=gray!10, right=6mm of sig] (aug) {Aumento:\\ninguno / clásico /\\difusión};
\node[box=blue!10, right=6mm of aug, yshift=7mm] (c1) {CNN 1D (WDCNN)};
\node[box=blue!10, below=1.5mm of c1] (c2) {STFT $\rightarrow$ CNN 2D};
\node[box=blue!10, below=1.5mm of c2] (c3) {16\,kHz $\rightarrow$ mel $\rightarrow$ AST};
\node[box=green!12, below=of c3] (out) {10 clases\\F1-macro, 3\,HP};
\draw[->] (sig) -- (aug); \draw[->] (sig) -- (ddpm); \draw[->] (ddpm) -| (aug);
\draw[->] (aug.east) -- (c1.west); \draw[->] (aug.east) -- (c2.west); \draw[->] (aug.east) -- (c3.west);
\draw[->] (c3) -- (out);
\end{tikzpicture}
\caption{Flujo del método: un único DDPM condicional genera ventanas de las clases escasas; tres clasificadores se entrenan con cada estrategia de aumento y se evalúan en la carga no vista.}
\label{fig:pipeline}
\end{figure}

\textbf{CNN 1D (línea base).} Arquitectura tipo WDCNN~\cite{zhang2017wdcnn}: una primera convolución de núcleo 64 y paso 16 seguida de cuatro bloques convolucionales (núcleo 3; BatchNorm, ReLU y \emph{max-pooling}), promedio global y un perceptrón de 100 unidades con \emph{dropout}~0,2.

\textbf{CNN 2D.} Sobre el espectrograma $S=\log(1+|\mathrm{STFT}(x)|)$ (ventana de Hann de 256 muestras, salto 8), normalizado por muestra; cuatro bloques convolucionales $3\times3$ (32--128 canales), promedio global y \emph{dropout}~0,3. Separa el efecto de la representación (1D frente a espectrograma) del de la arquitectura.

\textbf{AST.} Partimos del modelo \nolinkurl{MIT/ast-finetuned-audioset-10-10-0.4593}~\cite{gong2021ast}: el espectrograma mel de 128 bandas se divide en parches de $16\times16$ que se proyectan, se suman a embeddings posicionales y pasan por 12 capas de autoatención
\begin{equation}
\mathrm{Attn}(Q,K,V)=\mathrm{softmax}\!\left(QK^\top/\sqrt{d_k}\right)V.
\end{equation}
Como el AST se preentrenó con audio a 16~kHz, remuestreamos la ventana de 12 a 16~kHz antes del espectrograma (ventana de 25~ms y salto de 2~ms), de modo que cada banda mel corresponde a las mismas frecuencias en Hz que en el preentrenamiento; las bandas por encima de 6~kHz (Nyquist original) quedan prácticamente sin energía. Las 86 tramas resultantes se rellenan a 96 y los embeddings posicionales preentrenados ($12\times101$ parches) se interpolan bilinealmente a $12\times9$. Se reemplaza la cabeza por una capa lineal de 10 salidas y se ajusta toda la red.

\subsection{Aumento de datos}
\textbf{Sobremuestreo (control).} Se repiten ventanas reales de cada clase, elegidas al azar con reemplazo, hasta igualar la clase mayoritaria. Produce exactamente la misma cantidad de ventanas por clase que el aumento por difusión, pero sin información nueva; así, la diferencia entre ambos mide el aporte de las muestras generadas y no solo el de equilibrar las clases.

\textbf{Clásico.} Durante el entrenamiento, cada ventana recibe ruido blanco con SNR aleatoria entre 10 y 30~dB, un desplazamiento circular aleatorio y un escalamiento de amplitud en $[0{,}8;\,1{,}2]$; para los modelos sobre espectrograma se añade SpecAugment~\cite{park2019specaugment} (dos máscaras de frecuencia y de tiempo).

\textbf{Difusión condicional.} Entrenamos un DDPM~\cite{ho2020ddpm} sobre las ventanas de entrenamiento (desbalanceadas) con el proceso directo
\begin{equation}
q(x_t\mid x_0)=\mathcal{N}\!\left(\sqrt{\bar\alpha_t}\,x_0,\;(1-\bar\alpha_t)I\right),\quad \bar\alpha_t=\textstyle\prod_{s\le t}(1-\beta_s),
\end{equation}
con $T=1000$ y $\beta_t$ lineal entre $10^{-4}$ y $0{,}02$. Una U-Net 1D (canales 32--256, bloques residuales condicionados por FiLM~\cite{perez2018film} sobre la suma de los embeddings del paso $t$ y de la clase $y$) predice el ruido minimizando
\begin{equation}
\mathcal{L}=\mathbb{E}_{x_0,y,t,\epsilon}\left\|\epsilon-\epsilon_\theta\!\left(\sqrt{\bar\alpha_t}x_0+\sqrt{1-\bar\alpha_t}\,\epsilon,\,t,\,y\right)\right\|^2 .
\end{equation}
Con probabilidad 0,1 la etiqueta se reemplaza por una clase nula para aprender también el modelo incondicional, y en el muestreo se aplica guía sin clasificador~\cite{ho2022cfg}
\begin{equation}
\tilde\epsilon=(1+w)\,\epsilon_\theta(x_t,t,y)-w\,\epsilon_\theta(x_t,t,\varnothing),\quad w=2,
\end{equation}
con 50 pasos DDIM~\cite{song2021ddim} y pesos promediados por EMA. Se generan ventanas de cada clase hasta igualar la clase mayoritaria (\NSYN{} ventanas sintéticas), y estas se suman al entrenamiento. El generador nunca ve validación ni prueba.

\subsection{Entrenamiento y métricas}
Todos los clasificadores se entrenan con entropía cruzada, AdamW~\cite{loshchilov2019adamw} (decaimiento 0,01), calentamiento lineal del 5\,\% seguido de decaimiento coseno, recorte de gradiente en 5, precisión mixta y lotes de 64: 30 épocas con tasa $10^{-3}$ para las CNN y 8 épocas con $3\times10^{-5}$ para el AST. El DDPM se entrena 60 épocas con AdamW ($2\times10^{-4}$). Se conserva la época de mejor F1-macro en validación. La métrica principal es el F1-macro en 3~HP, que pondera igual a cada clase (la exactitud premiaría a la clase Normal); reportamos además el recall por clase, el F1-macro con ruido blanco añadido a 10 y 0~dB, y la media $\pm$ desviación estándar sobre las semillas 0, 1 y 2. La rejilla completa suma $3\times4\times3=36$ entrenamientos. Software: PyTorch \VTORCH{} y Transformers \VTRANSF{}; hardware: \HW.

% ===========================================================================
\section{Resultados}
\subsection{Línea base y ablaciones}
La Tabla~\ref{tab:resumen} compara las doce combinaciones. \pend{Describir: (1) cuánto mejora el mejor sistema sobre la CNN 1D sin aumento; (2) efecto de la representación (CNN 1D frente a CNN 2D); (3) efecto del Transformer preentrenado (CNN 2D frente a AST); (4) efecto del sobremuestreo, del aumento clásico y de la difusión en cada modelo, y en particular si la difusión supera al sobremuestreo, indicando si la diferencia supera la desviación entre semillas; (5) robustez al ruido.} La Fig.~\ref{fig:f1} muestra las mismas cifras.

\begin{table*}[t]
\centering
\caption{F1-macro en la carga no vista (3~HP), limpio y con ruido, media $\pm$ desviación sobre 3 semillas. En negrita, el mejor por columna.}
\label{tab:resumen}
\small
\tablagen{tabla_resumen}
\end{table*}

\begin{figure}[t]
\centering
\figgen{fig05_f1_por_modelo_aumento}{\columnwidth}
\caption{F1-macro en 3~HP por modelo y estrategia de aumento (barras: media; error: desviación entre semillas).}
\label{fig:f1}
\end{figure}

\textbf{Clases escasas.} La Tabla~\ref{tab:recall} muestra el recall de las clases más afectadas por el desbalance (severidades 0,014 y 0,021~in). \pend{Analizar si la difusión mejora específicamente las clases con 3\,\% de datos y cuáles se confunden (Fig.~\ref{fig:cm}).}

\begin{table}[t]
\centering
\caption{Recall de las clases minoritarias en 3~HP (semilla media).}
\label{tab:recall}
\resizebox{\columnwidth}{!}{\tablagen{tabla_recall}}
\end{table}

\begin{figure}[t]
\centering
\figgen{fig06_matrices_confusion}{\columnwidth}
\caption{Matrices de confusión en 3~HP.}
\label{fig:cm}
\end{figure}

\subsection{Calidad de las señales sintéticas}
\pend{Comentar la Fig.~\ref{fig:sintetico} (forma de onda y espectro real frente a sintético), la Fig.~\ref{fig:tsne} (t-SNE: ¿las sintéticas caen en los mismos grupos que las reales?) y la Tabla~\ref{tab:tstr} (TSTR: F1 de un clasificador entrenado solo con sintéticos y probado con reales).}

\begin{figure}[t]
\centering
\figgen{fig04_real_vs_sintetico}{\columnwidth}
\caption{Ventanas y espectros reales frente a sintéticos por clase.}
\label{fig:sintetico}
\end{figure}

\begin{figure}[t]
\centering
\figgen{fig11_tsne}{0.9\columnwidth}
\caption{t-SNE de los embeddings de la CNN 1D: ventanas reales y sintéticas.}
\label{fig:tsne}
\end{figure}

\begin{table}[t]
\centering
\caption{TSTR: entrenamiento solo con ventanas sintéticas, prueba con reales.}
\label{tab:tstr}
\resizebox{\columnwidth}{!}{\tablagen{tabla_tstr}}
\end{table}

\subsection{Interpretabilidad}\label{sec:interp}
Para cada ventana de prueba de las clases IR y OR calculamos el Grad-CAM~\cite{selvaraju2017gradcam} de la segunda etapa convolucional, cuya resolución temporal (375~Hz) permite resolver la periodicidad de BPFO y BPFI, y estimamos la frecuencia dominante del mapa de activación. La Tabla~\ref{tab:gradcam} compara esa frecuencia con las ecuaciones~\eqref{eq:bpfo}--\eqref{eq:bpfi}. \pend{Interpretar: porcentaje de ventanas cuya activación coincide (±12\,\%) con la frecuencia teórica; si la red atiende a los impactos del defecto o a otras componentes.} La Fig.~\ref{fig:ast} muestra la atención del AST (\emph{rollout}~\cite{abnar2020rollout}) por banda mel. \pend{Qué bandas atiende por clase.}

\begin{table}[t]
\centering
\caption{Frecuencia dominante del Grad-CAM frente a la frecuencia teórica de falla (3~HP).}
\label{tab:gradcam}
\tablagen{tabla_gradcam}
\end{table}

\begin{figure}[t]
\centering
\figgen{fig08_gradcam_2d}{\columnwidth}
\caption{Grad-CAM de la CNN 2D sobre el espectrograma para clases IR y OR.}
\label{fig:gradcam}
\end{figure}

\begin{figure}[t]
\centering
\figgen{fig10_ast_atencion}{\columnwidth}
\caption{Atención del AST por banda mel y clase.}
\label{fig:ast}
\end{figure}

\subsection{Costo computacional}
La Tabla~\ref{tab:costo} desglosa el tiempo de cómputo: el experimento completo requirió \GPUH~horas de GPU en \HW.

\begin{table}[t]
\centering
\caption{Costo computacional de la corrida completa.}
\label{tab:costo}
\footnotesize
\tablagen{tabla_costo}
\end{table}

% ===========================================================================
\section{Discusión}
\pend{Interpretación crítica de los resultados: por qué el aumento por difusión ayuda (o no) bajo cambio de carga; si el preentrenamiento en audio se transfiere a vibración; casos de error recurrentes (p.\,ej., confusión entre severidades de bola, que Smith y Randall~\cite{smith2015cwru} señalan como difíciles de diagnosticar).}

\textbf{Limitaciones.} (i) En CWRU cada condición de falla corresponde a un único rodamiento ensayado en las cuatro cargas; probar en 3~HP mide la generalización a una nueva condición de operación, no a un rodamiento nuevo, y puede sobrestimar el desempeño en campo. (ii) Los daños se introdujeron artificialmente por electroerosión; las fallas por desgaste real (Paderborn~\cite{lessmeier2016kat}) son más difusas. (iii) La clase Normal procede de una única grabación por carga. (iv) El AST se preentrenó con sonidos generales, no con vibración, y su entrada pierde la banda superior a 6~kHz. (v) Con 3 semillas, diferencias menores que la desviación no son concluyentes.

\textbf{Sesgos y uso responsable.} El modelo aprende un único banco de ensayo, motor y rodamiento; aplicarlo a otra máquina sin validación local puede producir falsas alarmas o, peor, fallas no detectadas. Recomendamos usarlo como apoyo al técnico de mantenimiento ---verificando con el espectro de envolvente que la alarma coincide con BPFO/BPFI--- y no como decisión automática de parada. Las señales sintéticas no deben usarse para validar ni reemplazar datos reales de prueba.

\textbf{Impacto.} El enfoque apunta a contextos con pocos registros de falla, como las pequeñas y medianas industrias peruanas, donde instalar un acelerómetro es barato pero acumular historiales de fallas toma años; la generación condicional permite entrenar con los pocos eventos disponibles.

% ===========================================================================
\section{Conclusiones}
\pend{2--4 hallazgos principales con números.} Como trabajo futuro proponemos validar en Paderborn y en datos de campo, filtrar las señales sintéticas con criterios físicos (presencia de BPFO/BPFI en su envolvente), y extender el enfoque a la estimación de vida útil remanente con redes informadas por física.

% ===========================================================================
\section*{Contribuciones}
% Ajustar con las tareas reales de cada integrante (el jurado puede preguntar a cada uno sobre lo que declara).
\begin{center}
\footnotesize
\begin{tabular}{p{0.3\columnwidth}p{0.6\columnwidth}}
\toprule
Integrante & Tareas en las que participó \\ \midrule
Jhafet Cánepa Maceda & Diseño experimental; implementación del paquete (modelos, DDPM, entrenamiento, interpretabilidad); repositorio \\
Lenin Dixon Herrera Padilla & Verificación y curación de datos (muestreo de la línea base, archivo duplicado); remuestreo del AST; reproducibilidad del repositorio; redacción del artículo \\
Hipólito Coaguila Berrios & \pend{completar: p.\,ej., ejecución de experimentos, análisis de resultados, diapositivas} \\
\bottomrule
\end{tabular}
\end{center}

\section*{Declaración de uso de asistentes de IA}
El equipo utilizó Claude (Anthropic) para revisar y depurar el código, verificar el conjunto de datos, documentar el repositorio y apoyar la redacción de este artículo. \pend{Añadir otras herramientas usadas por el equipo y para qué.} Todos los resultados provienen de ejecuciones propias del equipo y fueron revisados por sus integrantes.

\bibliographystyle{IEEEtran}
\bibliography{refs}

\end{document}
